\documentclass[10pt,a4paper]{amsart}
\usepackage[margin=19mm]{geometry}
\usepackage{amsmath,amssymb,mathtools}
\usepackage[scr=rsfs,cal=euler]{mathalfa}
\usepackage{xfrac}
\usepackage[unicode,hidelinks,bookmarks=false]{hyperref}
\newcommand{\ie}{i.e.\ }
\newcommand{\cf}{cf.\ }
\newcommand{\resp}{respectively\ }
\newcommand{\Subsection}{\subsection}
\providecommand{\nobreakdash}{\nobreak\hbox{-}}
\setlength{\emergencystretch}{2em}
\multlinegap0pt
\usepackage{mathtools}
\usepackage{amssymb}
\usepackage[shortlabels]{enumitem}
\newtheorem{proposition}{Proposition}
\newtheorem{lemma}{Lemma}
\usepackage{cleveref}
\crefname{proposition}{Proposition}{Propositions}
\crefname{lemma}{Lemma}{Lemmas}
\newcommand{\M}{\mathrm{m}}
\newcommand{\N}{\mathbb{N}}
\newcommand{\PS}[2]{X_{#2, #1}}
\newcommand{\RPA}[2]{#1 \cdot #2}
\newcommand{\composablepairs}[1]{#1^{(2)}}
\newcommand{\degree}{\mathbf{d}}
\newcommand{\morphdom}[1]{\mathrm{dom}(#1)}
\title{Groupoids of finitely aligned higher-rank graphs \\ via filters and graph morphisms}
\author{Lisa Orloff Clark, Malcolm Jones}
\date{Source: arXiv:2509.13929v2 (CC BY 4.0)}
\begin{document}
\maketitle

\noindent\emph{Source and licence.} Groupoids of finitely aligned higher-rank graphs \\ via filters and graph morphisms. Version arXiv:2509.13929v2 (CC BY 4.0), distributed by arXiv under the Creative Commons Attribution 4.0 International licence, http://creativecommons.org/licenses/by/4.0/, as recorded in the arXiv licence metadata for this exact version and shown on the abstract page https://arxiv.org/abs/2509.13929v2. Full licence notice: \texttt{licenses/arxiv-2509.13929-CC-BY-4.0.txt}.\par\smallskip

\noindent\textit{Editorial scope.} Counted unit: the article's lemma lem:action\_on\_PS\_of\_morphisms (source0.tex lines 795--805). The article's complete original proof of it is delivered with it (source0.tex lines 806--840, one \texttt{proof} environment of 35 lines). No mathematics is written, repaired or re-derived: the statement and the proof are the article's own lines, in the article's own notation, and no retained line is substituted or re-flowed.  Retained uncounted as the source-local context the proof needs: lines 713--718; lines 774--784.  Nothing was omitted from any retained span, so the package carries no omission record.  No retained line contains a citation, so the wrapper carries no bibliography and no bibliography entry.  No retained line contains a figure environment, an image-inclusion command or drawing code, so no image is delivered and the package declares no asset dependency.  Editorial formatting only, in addition: an AMS \texttt{amsart} preamble that restates the article's own declarations and macros used by the retained lines, namely the environments \texttt{proposition}, \texttt{lemma} on separate counters, together with the standard packages those lines need.  The retained environments print in this document as Proposition 1, Lemma 1 in order of appearance, whereas the article prints the same items with the numbering of its own full document.  The complete native source of the article as distributed in the arXiv e-print for this exact version is bundled unchanged as \texttt{sources/185202/source0.tex}.
\begin{proposition}
\label{prop:increasing_sequences_in_P_yield_nested_P_graphs}
Let $(Q, P)$ be a weakly quasi-lattice ordered group.
Suppose $(m_n) \subseteq P$ is $\leq$-increasing.
For all $n \in \N$, $\Omega_{P, m_n} \subseteq \Omega_{P, m_{n+1}}$ and the inclusion map from $\Omega_{P, m_n}$ to $\Omega_{P, m_{n+1}}$ is a $P$-graph morphism, and we denote the direct limit of $(\Omega_{P, m_n})$ by $\Omega_{P, (m_n)}$.
\end{proposition}

\begin{lemma}
\label{lem:actionable_morphisms}
Let $(Q, P)$ be a weakly quasi-lattice ordered group, and let $\Lambda$ be a finitely aligned $P$-graph. 
The following are equivalent:
\begin{enumerate}
\item \label{actionable_1} $(x, p) \in \PS{\M}{\Lambda} * P$;
\item \label{actionable_2} for all $\leq$-increasing $(m_n)$, $\morphdom{x} = \Omega_{P, (m_n)}$ implies $p \leq m_n$ eventually;
\item \label{actionable_3} there is some $\leq$-increasing $(m_n)$ with $\morphdom{x} = \Omega_{P, (m_n)}$ and $p \leq m_n$ eventually;
\item \label{actionable_4} there is some $\leq$-increasing $(m_n)$ with $\morphdom{x} = \Omega_{P, (m_n)}$ and $p \leq m_n$ for all $n$.
\end{enumerate}
\end{lemma}

\begin{lemma}
\label{lem:action_on_PS_of_morphisms}
Let $(Q, P)$ be a weakly quasi-lattice ordered group, and let $\Lambda$ be a finitely aligned $P$-graph. 
For each $(x, m) \in \PS{\M}{\Lambda} * P$, $\morphdom{x} = \Omega_{P, (m_n)}$ for some $\leq$-increasing $(m_n) \subseteq P$ with $m \leq m_n$ for all $n$.
The map
\[
\RPA{x}{m} \colon \Omega_{P, (m^{-1}m_n)} \to \Lambda, \quad 
(p, q) \mapsto (\RPA{x}{m})(p, q) \coloneqq x(mp, mq),
\]
is a graph morphism, and $(p, q) \in \morphdom{\RPA{x}{m}} \iff (mp, mq) \in \morphdom{x}$.
\end{lemma}

\begin{proof}
By \cref{lem:actionable_morphisms}, $\morphdom{x} = \Omega_{P, (m_n)}$ for some $\leq$-increasing $(m_n) \subseteq P$ with $m \leq m_n$ for all $n$.
By the definition of $\leq$ and because $\leq$ is left invariant, $(m^{-1}m_n)$ is a $\leq$-increasing subset of $P$, and hence has direct limit $\Omega_{P, (m^{-1}m_n)}$ as per \cref{prop:increasing_sequences_in_P_yield_nested_P_graphs}. 
Because $m \leq m_n$ for all $n$, $(p, q) \in \Omega_{P, (m^{-1}m_n)} \iff (mp, mq) \in \Omega_{P, (m_n)}$, which implies $\RPA{x}{m}$ is well-defined.

To see $\RPA{x}{m}$ is a functor, suppose $p \leq q \leq r \leq m^{-1}m_n$ for some $n$ (i.e. $((p, q), (q, r))$ is a composable pair in $\Omega_{P, (m^{-1}m_n)}$). 
By left invariance, 
\[
mp \leq mq \leq mr \leq m_n,
\] 
so $((mp, mq), (mq, mr)) \in \composablepairs{\Omega_{P, (m_n)}}$. 
Since $x$ is a functor, 
\begin{align*}
((\RPA{x}{m})(p, q), (\RPA{x}{m})(q, r)) 
&= (x(mp, mq), x(mq, mr)) \in \composablepairs{\Lambda} 
\text{ and } \\
(\RPA{x}{m})(p, q)(\RPA{x}{m})(q, r) 
&= x(mp, mq)x(mq, mr) \\ 
&= x((mp, mq)(mq, mr)) \\ 
&= x(mp, mr) = (\RPA{x}{m})(p, r) \\
&= (\RPA{x}{m})((p, q)(q, r)),
\end{align*} 
so $\RPA{x}{m}$ is a functor.
Also, 
\[
\degree((\RPA{x}{m})(p, q)) 
= \degree(x(mp, mq)) 
= \degree(mp, mq) 
= (mp)^{-1}mq 
= p^{-1}q 
= \degree(p, q),
\] 
so $\RPA{x}{m}$ is degree-preserving. 
That is, $\RPA{x}{m}$ is a graph morphism.
\end{proof}

\end{document}
